\section{Benchmark and Experimental Setup}
\label{sec:setup}

\subsection{Benchmark Task and Construction}

We construct a controlled benchmark for longitudinal proactive intervention with 207 sessions, 3,435 evaluated turns, and 301 positive intervention opportunities across six domains.
Each source-grounded record or scenario template is converted into the event representation from Section~\ref{sec:problem} and expanded into a multi-turn session in which a method must either Push or abstain at every turn.
Relative to ProactiveBench~\citep{lu2025proactiveagent}, which evaluates whether a generated proactive task would be accepted, our labels identify whether the current turn is an appropriate intervention opportunity independently of proposal wording.
The benchmark therefore isolates longitudinal evidence use, intervention timing, and abstention; it is not a reproduction of the ProactiveBench protocol.

\paragraph{Public-source grounding.}
Code uses a provenance pool of 989 non-merge commits collected from Requests (179), Flask (194), HTTPX (200), tqdm (216), and FastAPI (200).
Bug-fix, feature, documentation, configuration, and refactoring activity families from this pool inform 32 controlled repository workflows with edit, test, debug, documentation, and commit stages; the sessions are scenario expansions rather than direct commit replays.
SmartHome lifts 20 public SmartBench traces---15 anomaly and five normal traces---into Easy event sessions while preserving their routine or anomaly semantics; 12 controlled Hard variants test less explicit or exception-bearing cues (snapshot commit \texttt{98bec394})~\citep{horizonsinzqs2026smartbench}.
Email and Calendar are anchored to 25 distinct WorkBench email records and 25 distinct calendar events, respectively~\citep{styles2024workbench}.
Project Management uses 18 existing WorkBench task records plus four schema-compatible create-task requests for operations without a pre-existing target record.
Shopping uses the $\tau$-bench Retail task/tool schema and its product, order, and user resources to construct 32 sessions around six recurring preference or action patterns; it is a schema-grounded adaptation rather than a one-to-one replay of official trajectories~\citep{yao2025taubench}.

\paragraph{Longitudinal conversion and labels.}
The source-grounded records and scenario families are expanded into repeated evidence, an actionable or non-actionable current stage, and observable consequences.
Easy sessions expose clearer recurring evidence and intervention stages.
Each domain's 12 Hard sessions comprise four ambiguous late positives, four implicit positives, and four aligned or explicit-exception negatives.
Turn-level labels follow pre-specified domain construction rules that are independent of \system{}'s predictions and internal phase detector.
This controlled construction makes intervention timing reproducible, but the labels are not annotations of naturally occurring longitudinal agent--user interactions and the resulting scores are not official scores on the unmodified source benchmarks.

\subsection{Dataset Statistics}

Table~\ref{tab:data} reports only evaluated benchmark sessions and turns.

\begin{table}[t]
\centering
\caption{Benchmark statistics. Pos. denotes turns labeled for proactive intervention.}
\label{tab:data}
\small
\begin{tabular}{lrrrrr}
\toprule
& \multicolumn{2}{c}{Sessions} & & & \\
\cmidrule(lr){2-3}
Domain & Easy & Hard & Turns/session & Turns & Pos. \\
\midrule
Code & 20 & 12 & 12--17 & 448 & 51 \\
SmartHome & 20 & 12 & 12--17 & 440 & 42 \\
Shopping & 20 & 12 & 15--17 & 495 & 52 \\
Email & 25 & 12 & 16--20 & 684 & 52 \\
Calendar & 25 & 12 & 16--20 & 684 & 52 \\
Project Mgmt. & 25 & 12 & 16--20 & 684 & 52 \\
\midrule
Total & 135 & 72 & 12--20 & 3,435 & 301 \\
\bottomrule
\end{tabular}
\end{table}

The released artifact retains \texttt{clean} as the internal split identifier for Easy; this is a presentation alias and does not change any frozen data or result.

\subsection{Research Questions}

We study four questions: (RQ1) Does \system{} improve proactive decision quality over the evaluated memory proxies? (RQ2) Which mechanisms contribute under different domain structures? (RQ3) Are the effects stable across repetitions, protocols, and model configurations? (RQ4) What is the measured call and token composition of the framework?

\subsection{Baselines}

All methods receive the same domain data and current event, use the same ground-truth labels, and produce the same Push/no-Push output under protocol-aligned evaluators; \system{} and the baseline families are executed by separate implementations rather than a shared internal execution path.
\textbf{No Memory} uses only the current event and short local context.
\textbf{Retrieval-augmented generation (RAG)-only} stores raw historical events and retrieves similar observations.
\textbf{Summary-Reflection} compresses observations into a long-term natural-language summary and reflections.
\textbf{Generative Memory} follows a relevance--recency--importance retrieval scheme inspired by Generative Agents~\citep{park2023generative}.
These baselines are local proxy implementations constructed for this study, not official reproductions of the named systems.
Across all baseline agents, the decision prompt includes at most 10 recent events in addition to the current event.
RAG-only ranks raw events by Jaccard token overlap; RAG-only and Generative Memory cap retrieval at \texttt{top\_k}=8.
Summary-Reflection updates every 20 observations, flushes any nonempty remainder at the end of warm-up, and retains at most 12 reflections.
Generative Memory uses Jaccard overlap as relevance and scores observations as $0.55\,\mathrm{relevance}+0.25\exp(-\mathrm{age}/80)+0.20\,\mathrm{importance}$, where age is observation-index distance.
Importance starts at 0.45, adds 0.25 for predefined action verbs, 0.20 for predefined salient terms, and 0.10 for user-authored events, and is capped at 1.0.
The baselines intentionally exclude \system{}'s typed Habit Store, generalized preference abstraction, and consolidation mechanism.

\subsection{Evaluation Protocol}

The main \texttt{all} condition is a continual stateful Easy-to-Hard run: long-term method state persists across the subset boundary, while session-local cooldown is reset.
We separately audit split-reset and reversed-order protocols rather than treating \texttt{all} as a simple arithmetic union.
Before evaluated turns, Code receives 34 initialization events, SmartHome receives 50, and each remaining domain receives 54.
These warm-up events initialize longitudinal state, are excluded from Table~\ref{tab:data} and all reported metrics, and are provided to every stateful method where supported.
A per-memory cooldown of three events is likewise applied where supported to avoid giving \system{} an intervention-frequency advantage by protocol alone.

\subsection{Metrics and Statistical Tests}

We compute precision, recall, F1, accuracy, true positives (TP), false positives (FP), false negatives (FN), true negatives (TN), and Push count from complete decision traces, while the main table presents precision, recall, F1, FP, and Push count.
For ablations, 95\% confidence intervals (CIs) use 10,000 paired, stratified session-level bootstrap samples, preserving within-session dependence and Easy/Hard stratum sizes, with fixed seed \texttt{20260713}.
Because turns within a session are dependent, inferential claims rely on these paired session-level intervals rather than treating individual turns as independent observations.
For Code and Calendar, the bootstrap analyses use validated Full repeat-2 traces as references because the original Full logs predate structured tracing; the reference traces reproduce the frozen Full predictions and confusion matrices.

\subsection{Implementation and Reproducibility}

The main matrix uses DeepSeek-V4-Flash with temperature zero, domain-specific warm-up, LLM Think enabled, and cooldown three.
The package requires Python 3.10 or newer; the current bundled audit runtime reports Python 3.12.13, while exact historical environments are not recoverable for every run.
Package defaults are a 30-event discovery window, a 30-minute consequence window, high-confidence and forgetting thresholds of 0.8 and 0.2, a 10-event candidate cooldown, and a five-event feedback cooldown. The frozen evaluators use a 50-event discovery window and override the candidate cooldown to three events.
The instrumented HCPM ablation, repeatability, cross-model, and efficiency evaluators write one JSON Lines (JSONL) decision trace per evaluated turn, containing trigger hits, candidates, score components, phase, Think output, cooldown state, ground truth, and TP/FP/FN/TN classification.
The source-level audit validates 90/90 main-result artifacts (72 baseline JSON result files and 18 HCPM source logs), while the manifest separately validates 36,779/36,779 JSONL rows across 65 instrumented HCPM runs.
Application programming interface (API) keys and endpoint credentials are excluded from the artifact.
None of the six frozen benchmark evaluators supplies \texttt{query\_vector\_store}; consequently, the external-vector Semantic branch is inactive in all reported experiments.
The gains observed in the four domains with semantic surface variation---Shopping, Calendar, Email, and Project Management---come from generalized Habit Variant Match rather than Semantic vector retrieval.

